% !TeX root = ../../Book4.tex
% 纲要标题：### D.2 稳定模与同伦范畴
% 纲要细目（写作范围）：
% #### D.2.1 Frobenius 代数的有限维模
% #### D.2.2 双数代数的稳定模
% #### D.2.3 逐项分裂正合的复形

\section{稳定模与同伦范畴}
\label{b4:sec:D02}

稳定范畴的定义能用于相当不同的对象。一个模型从有限维代数出发：某些代数的投射模本身也内射。另一个模型从复形出发：若只指定逐项分裂的短正合列，可缩复形恰好承担投射与内射的共同角色。前者让我们计算一个具体模的稳定态射，后者则重新得到本卷已经熟悉的同伦范畴。

\subsection{Frobenius 代数的有限维模}
\label{b4:subsec:D02:01}

设 \(k\) 为域，\(A\) 为有限维含幺结合 \(k\)-代数，未必交换。记 \(A\text{-}\mathrm{mod}\) 为有限维幺性左 \(A\)-模及其 \(A\)-线性映射组成的范畴，配备通常的短正合结构。对偶空间
\[
D(A)=\operatorname{Hom}_k(A,k)
\]
有左 \(A\)-作用 \((a\varphi)(b)=\varphi(ba)\)。这里出现 \(ba\) 而非 \(ab\)，是为使所定义的作用成为左作用。
\symidx{\(A\text{-}\mathrm{mod}\)：有限维代数上的有限维左模范畴}

\begin{definition}\label{b4:def:D-frobenius-algebra}
设 \(k\) 为域，\(A\) 为有限维含幺结合 \(k\)-代数。在 \(D(A)=\operatorname{Hom}_k(A,k)\) 上指定左作用 \((a\varphi)(b)=\varphi(ba)\)。如果左正则模 \({}_AA\) 与 \({}_AD(A)\) 同构，就称 \(A\) 为 \term[frobenius-daishu]{Frobenius 代数}[Frobenius algebra]。
\end{definition}
\symidx{\(D(A)\)：有限维代数的线性对偶及其左模结构}

\begin{proposition}\label{b4:prop:D-frobenius-module-category}
设 \(k\) 为域，\(A\) 为 \(k\) 上的有限维 Frobenius 代数。则通常正合结构下的 \(A\text{-}\mathrm{mod}\) 是 Frobenius 范畴，投射内射对象恰为有限生成投射左 \(A\)-模。
\end{proposition}
\begin{proof}
对每个有限维左 \(A\)-模 \(M\)，有自然同构
\[
\operatorname{Hom}_A(M,D(A))\longrightarrow D(M),
\qquad f\longmapsto\bigl(m\mapsto f(m)(1)\bigr).
\]
其逆把 \(\lambda\in D(M)\) 送到
\[
f_\lambda(m)(b)=\lambda(bm).
\]
确有 \(f_\lambda(am)(b)=\lambda(bam)=(a f_\lambda(m))(b)\)，故逆映射左 \(A\)-线性。若 \(f\) 左线性，则
\(f(m)(b)=(b f(m))(1)=f(bm)(1)\)，所以两映射互逆。有限维向量空间的对偶函子正合，因此 \(D(A)\) 内射。假设 \(A\cong D(A)\) 便使左正则模 \(A\) 内射，有限直和及其直和项也内射；每个有限生成投射模是某个 \(A^r\) 的直和项，因而内射。

每个 \(M\) 接收有限自由模的满射，所以有足够多投射对象。另取 \(D(M)\) 的一组基 \(\lambda_1,\ldots,\lambda_d\)，以上公式给出单射
\[
M\longrightarrow D(A)^d,
\qquad m\longmapsto(f_{\lambda_1}(m),\ldots,f_{\lambda_d}(m)).
\]
若像为零，在 \(b=1\) 处求值便得所有 \(\lambda_j(m)=0\)，故 \(m=0\)。这证明有足够多内射对象。若 \(M\) 本身内射，上述单射分裂，于是 \(M\) 是 \(D(A)^d\cong A^d\) 的直和项，因而投射。两类对象相同，结论成立。
\end{proof}

这个证明只使用左模同构 \(A\cong D(A)\)，没有把 \(A\) 当作交换代数。稳定范畴 \(\underline{A\text{-}\mathrm{mod}}\) 因此把经过有限生成投射左模的态射视为零。

\subsection{双数代数的稳定模}
\label{b4:subsec:D02:02}

取任意域 \(k\) 上的双数代数
\[
A=k[\varepsilon]/(\varepsilon^2).
\]
令 \(\lambda(a+b\varepsilon)=b\)。映射
\[
\phi:A\longrightarrow D(A),\qquad
\phi(x)(y)=\lambda(yx)
\]
左 \(A\)-线性，因为 \(\phi(ax)(y)=\lambda(yax)=(a\phi(x))(y)\)。若 \(x=a+b\varepsilon\ne0\)，则 \(b\ne0\) 时 \(\lambda(x)=b\ne0\)，\(a\ne0\) 时 \(\lambda(\varepsilon x)=a\ne0\)。故 \(\phi\) 单射；两边维数均为二，所以同构。由\cref{b4:prop:D-frobenius-module-category}，\(A\text{-}\mathrm{mod}\) 是 Frobenius 范畴。

\begin{proposition}\label{b4:prop:D-dual-number-stable}
设 \(k\) 为域，\(A=k[\varepsilon]/(\varepsilon^2)\)，并把 \(k=A/(\varepsilon)\) 视为左 \(A\)-模。每个有限维左 \(A\)-模都同构于
\[
A^r\oplus k^s\qquad(r,s\ge0).
\]
其中投射模恰为 \(A^r\)，并且
\[
\operatorname{Hom}_{\underline{A\text{-}\mathrm{mod}}}(k^s,k^t)
\cong M_{t\times s}(k),
\qquad
\Omega k\cong k\cong\Sigma k.
\]
把有限维 \(k\)-向量空间赋予零 \(\varepsilon\) 作用，给出到 \(\underline{A\text{-}\mathrm{mod}}\) 的加性等价。
\end{proposition}
\begin{proof}
一个 \(A\)-模就是带有平方为零算子 \(T(m)=\varepsilon m\) 的向量空间。取 \(\operatorname{im}T\) 的基 \(u_1,\ldots,u_r\)，选 \(v_i\) 使 \(Tv_i=u_i\)，再把 \(u_i\) 扩充为 \(\ker T\) 的基 \(u_i,w_1,\ldots,w_s\)。这些 \(u_i,v_i,w_j\) 构成全空间的基：应用 \(T\) 可先消去任何线性关系中的 \(v_i\) 系数；每个向量减去适当的 \(v_i\) 组合后则落入 \(\ker T\)。每对 \(v_i,u_i\) 给出一份 \(A\)，每个 \(w_j\) 给出一份 \(k\)，得到分解。

模 \(k\) 不投射。若商映射 \(\pi:A\to k\) 有 \(A\)-线性截面 \(s\)，则 \(\varepsilon s(1)=s(\varepsilon\cdot1)=0\)，故 \(s(1)\in\varepsilon A\)，于是 \(\pi s(1)=0\)，与截面条件矛盾。因此含有 \(k\) 直和项的模不投射，而 \(A^r\) 自由，投射模恰为这些对象。

任意 \(A\)-线性映射 \(k^s\to A^r\) 的像落入 \(\varepsilon A^r\)，而任何 \(A^r\to k^t\) 都把该子模送到零。因此两者的复合为零。所有经过投射模的 \(k^s\to k^t\) 映射均为零，稳定 Hom 就等于普通 Hom，即所给矩阵空间。每个 \(A^r\oplus k^s\) 在稳定范畴中同构于 \(k^s\)，故向量空间的这个函子全忠实且本质满。具体的拟逆可取
\[
M\longmapsto\ker(\varepsilon:M\to M)/\varepsilon M;
\]
它把 \(A\) 送到零，把 \(k\) 送到 \(k\)，并由上面的分解给出同一稳定 Hom 识别。

最后，短正合列
\begin{equation}\label{b4:eq:D-dual-number-extension}
0\longrightarrow k\xrightarrow{\iota}A\xrightarrow{\pi}k\longrightarrow0,
\qquad \iota(1)=\varepsilon,
\end{equation}
既可作右端 \(k\) 的投射序列，也可作左端 \(k\) 的内射序列。于是 \(\Omega k=\varepsilon A\cong k\)，\(\Sigma k=A/\varepsilon A\cong k\)。标量映射可提升、延拓为 \(A\) 上的同一标量映射，所以这些识别也保持态射。在向量空间模型中，\(\Omega\) 与 \(\Sigma\) 因而都自然同构于恒等函子。
\end{proof}

上述等价是加性范畴的等价。原来的模范畴有非分裂扩张\eqref{b4:eq:D-dual-number-extension}，稳定商却把中项 \(A\) 变成零。在 \(\Sigma k\cong k\) 的选择下，这条扩张给出的好三角为
\[
k\longrightarrow0\longrightarrow k\xrightarrow{-1_k}k,
\]
连接态射的计算见\cref{b4:prop:D-conflation-triangle}后的双数代数例子。这并不是有限维向量空间范畴中的短正合列；加性等价本身并未将原来的模短正合结构搬到向量空间上。

\ProseParagraphBreak
双数代数中 \(\Omega\cong\operatorname{Id}\) 也不是 Frobenius 条件的一般后果。一般结论是 \(\Omega\) 与 \(\Sigma\) 互为拟逆；对三次截断多项式环的简单模，二次合冲才回到原模，一次合冲并不与原模稳定同构，具体计算见\cref{b4:prop:D-truncated-cubic-period}。

\subsection{逐项分裂正合的复形}
\label{b4:subsec:D02:03}

设 \(\mathcal A\) 为局部小加性范畴。上链复形与复形映射的定义仍可使用，因为只需态射加法、复合及零态射。给 \(\operatorname{Ch}(\mathcal A)\) 指定在每个次数都分裂的短正合列。分裂只针对各项，不要求所选截面组成复形映射。

\ProseParagraphBreak
例如，设 \(k\) 为域，\(D=(k\xrightarrow{1_k}k)\) 位于次数 \(0,1\)，\(S^i\) 是只在次数 \(i\) 取值 \(k\) 的复形。短正合列
\[
0\longrightarrow S^1\longrightarrow D\longrightarrow S^0\longrightarrow0
\]
逐项分裂，却不在复形范畴中分裂。若商映射有复形截面 \(s\)，零次截面条件要求 \(s^0=1_k\)，复形映射条件却要求
\[
\mathrm{d}_D^0s^0=s^1\mathrm{d}_{S^0}^0=0,
\]
而左边为 \(1_k\)。每度存在截面，只解决分次对象的分裂，尚未解决与微分的相容性。

\begin{proposition}\label{b4:prop:D-complex-frobenius}
设 \(\mathcal A\) 为局部小加性范畴。逐项分裂的短正合列给出 \(\operatorname{Ch}(\mathcal A)\) 的正合结构；在此结构下它是 Frobenius 范畴，投射内射对象恰为可缩复形。
\end{proposition}
\begin{proof}
对任意复形 \(X,Y\)，复形映射 \(X\to Y\) 由满足 \(\mathrm{d}_Y^nf^n=f^{n+1}\mathrm{d}_X^n\) 的态射族 \((f^n)_{n\in\mathbb Z}\) 给出。这些态射族构成逐项 Hom 集的可数积
\[
\prod_{n\in\mathbb Z}\operatorname{Hom}_{\mathcal A}(X^n,Y^n)
\]
的子集。因为 \(\mathcal A\) 局部小，这个可数积是集合，故 \(\operatorname{Ch}(\mathcal A)\) 也局部小。

逐项分裂列的核与余核在每个次数都有泛性质，所得态射族与微分相容，所以也在复形范畴中互为核与余核。选逐项分裂后，中项的微分有形式
\[
Y^n=X^n\oplus Z^n,
\qquad
\mathrm{d}_Y^n=
\begin{pmatrix}\mathrm{d}_X^n&t^n\\0&\mathrm{d}_Z^n\end{pmatrix},
\quad
\mathrm{d}_X^{n+1}t^n+t^{n+1}\mathrm{d}_Z^n=0.
\]
右上角 \(t^n\) 记录的正是逐项分裂后微分仍可能把商的坐标送入核的坐标；上面的两项复形例子中，它就是 \(1_k\)。
沿复形映射 \(f:X\to X'\) 推出，取 \(Y'^n=X'^n\oplus Z^n\)，把右上角改为 \(f^{n+1}t^n\)；这满足复形条件，逐项推出的泛性质也保证所得复形有推出泛性质。沿 \(g:Z'\to Z\) 拉回，则取 \(X^n\oplus Z'^n\)，右上角改为 \(t^ng^n\)。两个逐项分裂包含的复合逐项仍是分裂包含，商由两个互补项的直和逐项给出；两个逐项分裂投影的复合及其核有对偶构造。因此全部正合结构公理成立。

若 \(K\) 可缩，固定收缩 \(h_K\)，满足 \(\mathrm{d}_Kh_K+h_K\mathrm{d}_K=1_K\)。对任意复形 \(X\)，Hom 复形 \(\operatorname{Hom}^\bullet(K,X)\) 与 \(\operatorname{Hom}^\bullet(X,K)\) 都可缩。具体地，对次数 \(n\) 的齐次映射 \(a\)，两者的收缩分别为
\[
a\longmapsto(-1)^n a h_K,
\qquad
a\longmapsto h_Ka.
\]
代入 \(\partial a=\mathrm{d}a-(-1)^na\mathrm{d}\)，两次交叉项相消，得到 \(\partial H+H\partial=1\)。

给定逐项分裂满射 \(p:E\to Y\) 及复形映射 \(f:K\to Y\)，先逐项提升为次数零映射 \(a_0:K\to E\)，满足 \(pa_0=f\)。逐项提升未必与微分相容，其误差为 \(\partial a_0\)。因为 \(p\partial a_0=\partial f=0\)，误差经核 \(i:X\to E\) 唯一写成
\[
\partial a_0=i r,\qquad r\in\operatorname{Hom}^1(K,X).
\]
又有 \(i\partial r=\partial^2a_0=0\)，而 \(i\) 逐项单射，所以 \(\partial r=0\)。核复形中的误差因此是一个一次余圈。Hom 复形的收缩把它变成余边：令 \(b=H(r)\)，收缩等式给出 \(\partial b=r\)。修正逐项提升为
\[
a=a_0-i b,\qquad
\partial a=ir-i\partial b=0,\qquad pa=f.
\]
于是 \(a\) 是所需的复形映射，证明 \(K\) 投射。这里不需要 \(K\) 的各项在 \(\mathcal A\) 中投射；逐项分裂提供初始提升，整个可缩复形的同伦则消去其误差。

对偶地，给定逐项分裂单射 \(i:X\to E\) 及复形映射 \(f:X\to K\)，先逐项延拓为 \(a_0:E\to K\)。其误差满足 \((\partial a_0)i=0\)，所以经余核 \(q:E\to Y\) 写成 \(rq\)，其中 \(r\in\operatorname{Hom}^1(Y,K)\)。由 \(\partial^2a_0=0\) 及 \(q\) 逐项满射，仍有 \(\partial r=0\)。用 \(\operatorname{Hom}^\bullet(Y,K)\) 的收缩取 \(b=H(r)\)，则 \(a_0-bq\) 的微分为零，且在 \(X\) 上仍等于 \(f\)。这给出复形映射延拓，证明 \(K\) 内射。

对任意 \(X\)，令
\[
I_X=\operatorname{Cone}(1_X),\qquad
I_X^n=X^n\oplus X^{n+1},\qquad
\mathrm{d}_{I_X}^n=
\begin{pmatrix}\mathrm{d}_X^n&1\\0&-\mathrm{d}_X^{n+1}\end{pmatrix}.
\]
映射 \(h^n(y,x)=(0,y)\) 是收缩。有逐项分裂短正合列
\begin{equation}\label{b4:eq:D-contractible-embedding}
0\longrightarrow X\xrightarrow{x\mapsto(x,0)}I_X
\xrightarrow{(y,x)\mapsto x}X[1]\longrightarrow0.
\end{equation}
这给出足够多内射对象。将序列移位 \(-1\)，得到逐项分裂短正合列
\[
0\longrightarrow X[-1]\longrightarrow I_X[-1]\longrightarrow X\longrightarrow0.
\]
\(I_X[-1]\) 仍可缩，因而由已证部分投射；到 \(X\) 的可容许满射给出足够多投射对象。任何内射对象都是相应 \(I_X\) 的直和项，任何投射对象都是相应 \(I_X[-1]\) 的直和项；若 \(ba=1_X\) 且 \(K\) 的收缩为 \(h\)，则 \(bha\) 收缩 \(X\)。因此这些直和项都可缩，两类对象恰为可缩复形。
\end{proof}

\begin{theorem}\label{b4:thm:D-stable-complex-homotopy}
设 \(\mathcal A\) 为局部小加性范畴，\(\operatorname{Ch}(\mathcal A)\) 配备逐项分裂正合结构。复形映射 \(f:X\to Y\) 通过可缩复形分解，当且仅当 \(f\) 零同伦。因此恒等对象对应诱导加性同构
\[
\underline{\operatorname{Ch}(\mathcal A)}\cong K(\mathcal A).
\]
在式\eqref{b4:eq:D-contractible-embedding}的选择下，稳定范畴的 \(\Sigma\) 就是复形移位 \([1]\)。
\end{theorem}
\begin{proof}
若 \(f=ba\) 通过可缩复形 \(K\)，取其收缩 \(h_K\)，则
\[
f=b(\mathrm{d}_Kh_K+h_K\mathrm{d}_K)a
=\mathrm{d}_Y(bh_Ka)+(bh_Ka)\mathrm{d}_X,
\]
所以零同伦。反过来，若 \(f=\mathrm{d}_Yh+h\mathrm{d}_X\)，定义
\[
b:I_X\longrightarrow Y,
\qquad b^n(y,x)=f^n(y)+h^{n+1}(x).
\]
复形条件在第二个分量上恰好为
\(\mathrm{d}_Y^nh^{n+1}=f^{n+1}-h^{n+2}\mathrm{d}_X^{n+1}\)，由零同伦等式成立；第一个分量则由 \(f\) 为复形映射保证。于是 \(b\) 是复形映射，并且 \(f=b\circ i_X\)。这给出经过可缩 \(I_X\) 的分解。

两个复形映射的稳定类相同，等价于它们之差零同伦，故商 Hom 群与 \(K(\mathcal A)\) 的 Hom 群完全相同，且保持复合。序列\eqref{b4:eq:D-contractible-embedding}的余核是 \(X[1]\)；对 \(f\) 的延拓可取 \((y,x)\mapsto(fy,fx)\)，诱导的余核映射正是 \(f[1]\)。因此移位也相同。
\end{proof}

这个识别使用逐项分裂的正合结构，对任意局部小加性范畴 \(\mathcal A\) 都成立，也不要求它幂等完备。若 \(\mathcal A\) 进一步是交换范畴，则可讨论零调复形与拟同构。在 \(K(\mathcal A)\) 中将拟同构变成同构，等价于把零调复形变成零，所得导出范畴为
\[
D(\mathcal A)=K(\mathcal A)[\mathrm{qis}^{-1}],
\]
其中 \(\mathrm{qis}\) 表示拟同构组成的态射类。稳定商消去可缩复形，导出局部化还会消去不可缩的零调复形。例如，对域 \(k\) 上的 \(A=k[\varepsilon]/(\varepsilon^2)\)，令 \(E\) 每个整数次数都为 \(A\)，每个微分均为乘 \(\varepsilon\)。\cref{b4:exm:unbounded-free-not-k-projective}已证明它零调但不可缩。因此 \(E\) 在同伦范畴及相应的逐项分裂稳定范畴中不为零，而 \(E\to0\) 是拟同构，故它在导出范畴中为零。复形的这一 Frobenius 模型也见 Bühler\cite[Section 13.4, Example 13.8]{Buehler2010ExactCategories}。

\begin{proposition}\label{b4:prop:D-truncated-cubic-period}
设 \(k\) 为域，\(A=k[t]/(t^3)\)，\(S=A/(t)\)，并在有限维幺性左 \(A\)-模范畴中取通常短正合结构。则 \(A\) 为 Frobenius 代数。在由商映射 \(A\to S\) 及乘 \(t\) 的满射 \(A\to tA\) 所作的合冲选择下，
\[
\Omega S=tA,\qquad \Omega^2S=t^2A\cong S,
\qquad
\operatorname{End}_{\underline{A\text{-}\mathrm{mod}}}(S)\cong k.
\]
但 \(\Omega S\) 与 \(S\) 在稳定范畴中不同构，所以 \(S\) 的合冲周期恰为二。
\end{proposition}
\begin{proof}
令 \(\lambda:A\to k\) 取 \(t^2\) 的系数，定义
\[
\phi:A\longrightarrow D(A),\qquad \phi(a)(b)=\lambda(ba).
\]
按 \(D(A)\) 的左作用，有 \(\phi(ca)(b)=\lambda(bca)=(c\phi(a))(b)\)，所以 \(\phi\) 左 \(A\)-线性。若非零 \(a\) 的最低非零次数为 \(j\in\{0,1,2\}\)，取 \(b=t^{2-j}\)，则 \(\lambda(ba)\ne0\)，故 \(\phi\) 单射。两边的 \(k\) 维数都是三，所以它同构，证明 Frobenius 性。由\cref{b4:prop:D-frobenius-module-category}，模范畴因此是 Frobenius 范畴，有限生成投射模正是所消去的投射内射对象。

商映射 \(A\to S\) 的核为 \(tA\)，而乘 \(t\) 的满射 \(A\to tA\) 的核为 \(t^2A\)。两张满射的来源都是投射模，故给出所列合冲。映射 \(S\to t^2A\)、\(\overline1\mapsto t^2\) 是左模同构，所以二次合冲回到 \(S\)。

普通自同态环 \(\operatorname{End}_A(S)\) 为 \(k\)。任意 \(S\to A^r\) 的像被 \(t\) 零化，因而落入 \(t^2A^r\)，而任何 \(A^r\to S\) 都杀掉该子模。所以经过有限自由模的 \(S\to S\) 映射为零。有限生成投射模是某个 \(A^r\) 的直和项，经过它的分解也可写成经过 \(A^r\) 的分解。因此商去投射分解后，自同态环仍为 \(k\)，特别有 \([1_S]\ne0\)。

任意 \(u:S\to tA\) 的像都被 \(t\) 零化，所以包含于 \(t^2A\)。任意 \(v:tA\to S\) 又满足 \(v(t^2)=t\cdot v(t)=0\)，故 \(vu=0\)。若 \(tA\) 与 \(S\) 稳定同构，就能选取这样两张映射，使其复合的稳定类为 \([1_S]\)，与刚证的 \([1_S]\ne0\) 矛盾。因此一次合冲未回到原对象，而二次合冲已回到它。
\end{proof}
